\documentclass[11pt,a4paper]{article}
\usepackage[spanish,es-nodecimaldot]{babel}
\usepackage[T1]{fontenc}
\usepackage[utf8]{inputenc}
\usepackage{geometry,graphicx,booktabs,tabularx,hyperref,xcolor}
\geometry{margin=2.3cm}
\hypersetup{colorlinks=true,urlcolor=blue!55!black}
\setlength{\parskip}{.55em}
\setlength{\parindent}{0pt}
\begin{document}
\pagenumbering{gobble}
\begin{titlepage}
\centering
\vspace*{1cm}
\includegraphics[width=5.3cm]{recursos/ufro.png}\par
\vspace{1cm}
{\Large UNIVERSIDAD DE LA FRONTERA\par}
{\large Ingeniería Informática\par}
\vspace{2.2cm}
{\Huge\bfseries Informe ejecutivo\par}
\vspace{.5cm}
{\Large Infraestructura de servicios para Fauna Araucanía\par}
{\large Taller de Redes · Grupo 3\par}
\vfill
\begin{tabular}{rl}
Integrantes: & Luis Jaramillo, Giovanny Toledo y Maximiliano Rivas\\
Profesor: & Flavio Rojas Acuña\\
Fecha de revisión: & 27 de septiembre de 2026\\
\end{tabular}
\vspace{1.3cm}
\end{titlepage}
\pagenumbering{arabic}

\section*{Necesidad y solución}
Fauna Araucanía necesitaba presencia en Internet y servicios básicos alojados en el VPS asignado. El Grupo 3 organizó la solución bajo \texttt{zorro-darwin.lazos.cl}: un DNS propio, tres sitios independientes, administración de DNS, transferencia de archivos, correo con acceso web y un programa de chat. El VPS usa AlmaLinux 9.8; los servicios se mantienen bajo firewall y SELinux en modo \texttt{Enforcing}.

\section*{Resultados comprobados}
\begin{tabularx}{\linewidth}{@{}p{3cm}X@{}}\toprule
Área & Resultado observado desde el VPS o un equipo externo \\\midrule
DNS y web & El dominio raíz y los nombres de servicio resolvieron a 200.13.5.39. WordPress, Joomla y Grav mostraron contenido propio en web1, web2 y web3. \\
Administración & PowerAdmin permitió ingresar y administrar la zona; los cambios de prueba fueron visibles en consultas DNS. \\
Archivos & Cada uno de los tres usuarios FTP pudo subir y modificar su archivo; la versión nueva se descargó desde su sitio. No pudieron entrar en carpetas de otros usuarios; el acceso anónimo fue rechazado. \\
Correo & CMS1 envió un mensaje desde RoundCube a CMS2; CMS2 lo abrió y respondió; CMS1 abrió la respuesta. \\
Chat y protección & El cliente local intercambió mensajes en dos conexiones consecutivas. Nueve servicios estaban activos; MariaDB solo escuchaba localmente y SELinux seguía en Enforcing. \\
\bottomrule\end{tabularx}

\section*{Dificultades y decisiones}
La revisión de la primera entrega detectó fallas que no aparecían en comprobaciones superficiales: FTP aceptaba usuarios pero rechazaba la escritura, y RoundCube permitía entrar pero no enviar. Se aplicó una regla SELinux acotada para la escritura FTP y se corrigió el destino SMTP de RoundCube. Se completó el contenido propio de Joomla y Grav, se añadió la dirección del dominio raíz y se retiraron dos artefactos temporales de prueba. Cada corrección se contrastó con una operación completa; el informe técnico contiene comandos, capturas explicadas y resultados.

\section*{Estado y límites}
La infraestructura funciona para las actividades comprobadas del taller. El webmail y el FTP públicos usan transporte sin cifrado, y el certificado de correo es autofirmado: antes de un uso real deben incorporarse HTTPS, un método seguro de transferencia y un certificado confiable. La cuenta adicional \texttt{admin@} mencionada en la planificación no fue demostrada. El video resumen solicitado por la pauta no se elaboró en esta revisión. El informe técnico es el respaldo detallado de esta síntesis; las credenciales no aparecen en ninguno de los dos documentos.

\end{document}
